% ============================================
% Proof of Theorem 1: WAIT Algorithm Optimality
% ============================================

\section{Proof of Theorem~\ref{thm:wait_heavy_traffic}}\label{appexidx::proof of wait_heavy_traffic}

\paragraph{Proof Outline.}
This proof establishes asymptotic optimality of the WAIT algorithm when output lengths are known. The structure proceeds in two parts:
\begin{enumerate}
    \item \textbf{Single-type case}: We construct a coupled dominating process using the Lindley recursion and apply Kingman's bound for queue length expectations.
    \item \textbf{Multiple-type case}: We introduce a conservative periodic comparison policy whose slots have the full-threshold processing time. This separates the random batch composition of event-driven WAIT from the deterministic clock used in the drift argument.
\end{enumerate}
Key techniques include: Lindley recursion for queue length bounds \citep[Proposition 6.3]{asmussen2003applied}, coupling construction to dominate the original process, Kingman's bound for negative drift processes \citep{kingman1962queues}, and standard random walk results \citep{strait1974maximum}.

Throughout this appendix, we use $b$ to denote the embedded comparison-slot index, which is distinct from the stage index $s \in \{0, 1, \ldots, l_j'\}$ used in the main text for decode stages. We use $t$ for continuous time and \(B\) for the number of comparison slots.

We first consider the single-type case, then extend to multiple types.

\subsubsection{Single-Type Case}

Consider a single prompt type with threshold $n$. For general scheduling algorithms, analyzing this system in continuous time is challenging: KV cache grows dynamically during decode, arrivals are stochastic, and eviction risks complicate the dynamics.

The WAIT policy's \emph{threshold mechanism achieves a dimensionality reduction}. By waiting until exactly $n$ prompts accumulate before initiating a batch, the algorithm enforces a fixed batch size, which ensures constant processing time $\Delta T$ per batch. This decouples the continuous-time memory-constrained dynamics into a tractable discrete-time process: a memory-constrained queueing system becomes a simple random walk. The reduction is enabled by the algorithm design and is what makes the subsequent drift analysis apply.

Specifically, let \(\tilde{\lambda}\) denote the continuous-time arrival rate and let \(\mu=\tilde{\lambda}\Delta T\) be the expected number of arrivals during a threshold-sized batch. The boundary case is \(\mu=n\). In this case the embedded process observed at batch-completion epochs, indexed by \(b\in\mathbb{N}\), is a zero-drift threshold queue. The throughput loss arises from ``stuck'' iterations where insufficient prompts are available.

\begin{lemma}[Queue Length and Stuck Time]\label{lemma::throughput gap, single}
This lemma connects queue accumulation to throughput loss through stuck iterations in the critical case \(\mu=n\). Let \(X^b\) denote the number of arrivals during batch \(b\), where \(X^b \sim \mathrm{Poisson}(n)\). Let \(W^b\) denote the queue length after batch \(b\), with \(W^0=0\). The state transition is:
$$W^{b+1}=W^b+X^b-n\cdot \mathbf{1}\{W^b+X^b\geq n\}.$$
Define \(B_{\emph{stuck}}\) as the number of stuck iterations, where \(W^b+X^b<n\). Then:
$$\mathbb{E}[W^B]= n\cdot\mathbb{E}[B_{\emph{stuck}}].$$
\end{lemma}
\noindent The proof is in Appendix~\ref{appendix_lemma::proof of throughput gap, single}. This lemma connects queue accumulation to throughput loss: stuck iterations directly translate to waiting prompts. To bound $\mathbb{E}[W^B]$, we construct a coupled dominating process that admits tractable analysis via classical queueing techniques.

\begin{lemma}[Coupled Dominating Process]\label{lem:coupled_process, single}
The coupled process starts with a safety buffer (\(2n\) versus 0) and maintains dominance throughout. Define \(\{\tilde{W}^b\}\) by:
$$
\tilde{W}^0 = 2n, \quad \tilde{W}^{b+1} = \max\{ 2n, \tilde{W}^b + X^b - n \}.
$$
Then \(\tilde{W}^b \geq W^b + n\) for all \(b\in \mathbb{N}\).
\end{lemma}
\noindent The proof is in Appendix~\ref{appendix_lemma::proof of coupled_process, single}.

\paragraph{Lindley recursion.} To analyze the coupled process, we apply the Lindley recursion representation \citep[Proposition 6.3]{asmussen2003applied}:
\begin{lemma}[Lindley Representation]\label{lem:lindley_rep}
Let \(S^k=\sum_{r=1}^k (X^r-n)\) be the partial sum process. Then:
$$\tilde{W}^B=2n+\max(S^B, S^B-S^1,\ldots,S^B-S^{B-1},0).$$
\end{lemma}
\noindent This classical queueing representation is made applicable by the WAIT policy's threshold design. Define \(\xi^k = X^k-n\). We use time-reversal symmetry to simplify the maximum expression:
$$\max(S^B, S^B-S^1,\ldots,S^B-S^{B-1},0)=\max_{1\leq k\leq B}\Big(\sum_{r=k}^{B}\xi^r\Big)^+\overset{d}{=}\max_{1\leq k\leq B}\Big(\sum_{r=1}^{k}\eta^r\Big)^+, \quad \eta^r\overset{d}{=}X^1-n.$$
Thus \(\tilde{W}^B \overset{d}{=}2n+\max_{1\leq k\leq B}(\sum_{r=1}^{k}\eta^r)^+\).

Since \(\tilde{W}^B \geq W^B + n\), we have \(\mathbb{E}[W^B]\leq \mathbb{E}[\tilde{W}^B]-n\). Let \(Y^B = \max_{1\leq k\leq B}(\sum_{r=1}^{k}\eta^r)^+\). By \cite[Lemma 2]{strait1974maximum}:
\begin{lemma}[Random Walk Maximum]\label{lem:Y_T}
$$\ex{}{Y^B}=\ex{}{\max_{1\leq k\leq B}(S^k)^+}=\sum_{k=1}^B \frac{1}{k}\ex{}{(S^k)^+}.$$
\end{lemma}
\noindent This standard result \citep{strait1974maximum} decomposes the maximum into a sum over batch indices. For \(\mathbb{E}[(S^k)^+]\), the Cauchy-Schwarz inequality and \(\mathrm{Var}(S^k)=nk\) imply \(\mathbb{E}[(S^k)^+] \leq \sqrt{nk}\). Hence:
$$\mathbb{E}[Y^B] \leq \sqrt{n} \sum_{k=1}^B k^{-1/2}=O(\sqrt{nB}).$$

\begin{lemma}[Throughput Gap Bound]\label{lemma::throughput gap}
$$\mathbb{E}[W^B]-n\leq \mathbb{E}[Y^B]\leq  \sqrt{n} \sum_{k=1}^B \frac{1}{\sqrt{k}}=O(\sqrt{nB}).$$
\end{lemma}
\noindent By Lemma~\ref{lemma::throughput gap, single}, the expected number of stuck iterations is \(\mathbb{E}[B_{\emph{stuck}}]=\mathbb{E}[W^B]/n=O(\sqrt{B})\), so the normalized throughput loss from stuck iterations is \(O(B^{-1/2})\). The same bound applies to every prefix \(b\le B\), and therefore
\[
    \frac{1}{B}\sum_{b=1}^B \mathbb{E}[W^b] \leq \frac{1}{B}\sum_{b=1}^B O(\sqrt{b}) = O(\sqrt{B}).
\]
This time-averaged backlog bound is the quantity used for the latency and TTFT estimates.

\paragraph{Scaling with the horizon.} Let \(\Delta^\pi\) denote the fixed processing time of a threshold-sized batch under policy \(\pi\). In the scaled system, the corresponding processing time is \(\Delta^\pi/\zeta\). Over a calendar horizon \(T\), set
\[
    B_\zeta=\left\lfloor \frac{\zeta T}{\Delta^\pi}\right\rfloor .
\]
Then \(B_\zeta=\zeta T/\Delta^\pi+O(1)\). The terminal backlog bound gives normalized throughput loss \(O(B_\zeta^{-1/2})=O((\zeta T)^{-1/2})\). For delay, the theorem reports service-normalized time, so one scaled comparison slot has length \(\Delta^\pi\); the time-averaged backlog bound therefore gives latency and TTFT contributions \(O(B_\zeta^{1/2})=O((\zeta T)^{1/2})\). The bounded number of service stages is lower order. This yields the single-type throughput, latency, and TTFT orders stated in the theorem.

\subsubsection{Multiple-Type Case}

We now extend to \(m\) known prompt types. The additional issue is that event-driven WAIT need not process the full type composition in every realized batch: if only a subset of type queues has reached its threshold, only that subset is batched. The proof therefore separates two objects. The actual policy is the event-driven WAIT policy from the main text. The comparison object is a periodic policy with deterministic full-threshold slots; it is slower than actual WAIT because it reviews less often and pads every selected batch to the full-slot length.

For a threshold vector \(\pi=[n_1,\ldots,n_m]\), let
\begin{equation}\label{eq:batch_time_def}
\Delta^\pi \equiv \Delta T_{[1,\ldots,m]}
    =  d_0 + d_1 M^\pi,\qquad
M^\pi = \sum_{j=1}^m n_j(l_j^\prime+1)\left( l_j+\frac{l_j^\prime}{2} \right).
\end{equation}
This is the processing time of the full threshold composition, i.e., the batch that contains \(n_j\) prompts from every decode stage of every type \(j\). The load-balance and memory conditions are
\begin{equation}\label{eq:policy_constraints}
\begin{aligned}
&\lambda_j\Delta^\pi \leq n_j \quad \text{for all } j \in [m], \\
&M^{*}\leq M^{\pi}\leq C.
\end{aligned}
\end{equation}
The first line says that, over one full-threshold iteration, the expected type-\(j\) arrivals do not exceed the \(n_j\) type-\(j\) prompts that can be advanced from each active stage. The second line is the corresponding feasibility condition: the threshold composition is at least on the scale of the ideal-fluid memory requirement but does not exceed the physical memory \(C\).

\begin{lemma}[Resident-stage memory invariant]\label{lem:wait_resident_memory_invariant}
Starting from the empty system, WAIT maintains
\[
    N_{js}(t)\le n_j,\qquad j\in[m],\; s=1,\ldots,l_j^\prime,
\]
where \(N_{js}(t)\) is the number of resident type-\(j\) prompts at decode stage \(s\). Stage-0 prompts that have not been selected for prefill have no resident KV cache.
\end{lemma}
\noindent The proof is by induction over batch completions. If type \(j\) is not served in a realized batch, its resident stages are unchanged. If type \(j\) is served, the invariant before the batch implies that each resident stage contains at most \(n_j\) prompts, so WAIT selects all prompts currently present in that stage. After the batch completes, stage \(s+1\) contains only the prompts advanced from stage \(s\), again at most \(n_j\). Thus no resident stage can exceed \(n_j\).

The invariant gives the memory-feasibility statement needed by the policy. If \(a_{j0}(t)\le n_j\) denotes the number of type-\(j\) prompts selected for prefill in the current batch, then at any decision epoch,
\[
    \sum_{j=1}^m\sum_{s=1}^{l_j^\prime}N_{js}(t)(l_j+s)
    +
    \sum_{j=1}^m a_{j0}(t)l_j
    \le
    \sum_{j=1}^m n_j\sum_{s=0}^{l_j^\prime}(l_j+s)
    =
    M^\pi
    \le C .
\]
Thus unselected resident prompts are included in the memory accounting, while unselected stage-0 prompts remain outside the GPU KV-cache state.

For any realized event-driven WAIT batch \(r\), let \(a^{(r)}_{js}\le n_j\) be the number of selected type-\(j\) prompts at decode stage \(s\), and let \(J_r\) be the set of selected types. The selected batch composition has size
\[
    M_r
    =
    \sum_{j\in J_r}\sum_{s=0}^{l_j^\prime} a^{(r)}_{js}(l_j+s)
    \le
    \sum_{j=1}^m n_j\sum_{s=0}^{l_j^\prime}(l_j+s)
    =
    M^\pi.
\]
Therefore the physical processing time of the realized batch in the \(\zeta\)-scaled system is
\[
    S_r^{(\zeta)}
    =
    \frac{d_0+d_1M_r}{\zeta}
    \le
    \frac{\Delta^\pi}{\zeta}.
\]
This bound concerns processing time, while Lemma~\ref{lem:wait_resident_memory_invariant} gives memory feasibility. Actual batches may contain fewer types and have shorter duration, but no realized WAIT batch is larger than the full-threshold composition.

In the \(\zeta\)-scaled system, one full-threshold slot has physical length \(\Delta^\pi/\zeta\). Define deterministic review times
\[
    u_b=b\Delta^\pi/\zeta,\qquad b=0,1,2,\ldots .
\]
The comparison policy \(\bar\pi\) reviews the entry queues only at these times. We use an end-of-slot accounting convention: arrivals accumulate over \((u_b,u_{b+1}]\), and at \(u_{b+1}\) the comparison policy removes one type-\(j\) entry cohort of size \(n_j\) whenever the type-\(j\) entry queue has reached \(n_j\); it also advances up to \(n_j\) resident type-\(j\) prompts from each active decode stage. Equivalently, one can view the selected work as occupying the following full slot. A cohort removed at review \(u_{b+1}\) can therefore shift completion accounting by at most \(l_j^\prime+1\) comparison slots for type \(j\), which contributes only \(O((\zeta T)^{-1})\) to normalized throughput and \(O(1)\) to service-normalized average delay. Every selected batch is padded to length \(\Delta^\pi/\zeta\), and the construction is feasible by the resident-stage invariant above.

\begin{lemma}[One-slot delayed comparison dominance]\label{lem:wait_one_slot_dominance}
Couple actual event-driven WAIT and the periodic comparison policy \(\bar\pi\) on the same arrival sample path, both starting empty. Let \(L_\zeta=\Delta^\pi/\zeta\) and \(u_b=bL_\zeta\). For type \(j\), let \(R_j^\pi(t)\) be the cumulative number of type-\(j\) threshold cohorts of size \(n_j\) started by actual WAIT by time \(t\), and let \(R_j^{\bar\pi}(b)\) be the cumulative number of type-\(j\) cohorts removed from the comparison entry queue at reviews \(u_1,\ldots,u_b\). Then
\[
    R_j^\pi(u_{b+1})\ge R_j^{\bar\pi}(b),
    \qquad b\ge0,\; j\in[m].
\]
The same one-slot delayed dominance holds for type-\(j\) stage advancements.
\end{lemma}
\noindent To see this, use induction over \(b\). If \(\bar\pi\) removes a type-\(j\) cohort at \(u_b\), then that cohort has accumulated by \(u_b\). Actual WAIT reviews the system at the arrival that first makes the cohort eligible and at subsequent batch-completion epochs. If actual WAIT has not already started the corresponding type-\(j\) cohort by \(u_b\), then it must be processing an earlier realized batch. That batch finishes by \(u_b+L_\zeta=u_{b+1}\), because every realized batch has duration at most \(L_\zeta\). Hence actual WAIT starts the cohort by \(u_{b+1}\). For resident stages, each type-\(j\) service opportunity advances at most one threshold cohort from every stage; applying the same induction over the service-opportunity count and then over the stage index gives the stated one-slot dominance for stage advancements.

Lemma~\ref{lem:wait_one_slot_dominance} means that the comparison policy is a delayed and padded version of actual WAIT. It is therefore enough to bound \(\bar\pi\): the one-slot lag and the bounded in-service pipeline contribute only \(O((\zeta T)^{-1})\) to normalized throughput and \(O(1)\) to the service-normalized delay metrics.

For the comparison policy, let \(\bar W^b_{(j)}\) be the residual type-\(j\) entry queue after the \(b\)-th review, and let
\[
    \bar X^b_{(j)}
    =A_j(u_b,u_{b+1}]
    \sim \text{Poisson}\!\left(\zeta\lambda_j\frac{\Delta^\pi}{\zeta}\right)
    =\text{Poisson}(\lambda_j\Delta^\pi).
\]
The arrivals \(\{\bar X^b_{(j)}\}\) are independent across comparison slots and do not depend on other type queues. With the end-of-slot accounting convention above, the residual entry queue evolves as
\[
    \bar W^{b+1}_{(j)}
    =
    \bar W^b_{(j)}+\bar X^b_{(j)}
    -
    n_j\,\mathbf{1}\{\bar W^b_{(j)}+\bar X^b_{(j)}\ge n_j\}.
\]
This is the point at which the batch-composition issue disappears: the only coupling across types is the deterministic feasibility condition \(M^\pi\le C\); under the comparison clock, each type has its own threshold queue.

\begin{lemma}[Multiple-Type Comparison Queue Bound]\label{lemma::multiple-type coupled dominace}
For each type \(j\), define
\[
\widetilde W_{(j)}^0=2n_j,\qquad
\widetilde W_{(j)}^{b+1}
=\max\{2n_j,\widetilde W_{(j)}^b+\bar X^b_{(j)}-n_j\}.
\]
Then \(\widetilde W^b_{(j)}\ge \bar W^b_{(j)}\) for all \(b\ge0\) and \(j\in[m]\).
\end{lemma}
\noindent The proof is in Appendix~\ref{appendix_lemma::proof of multiple-type coupled dominace}. Let \(\mu_j=\lambda_j\Delta^\pi\). The increment in the dominating process is \(\bar X^b_{(j)}-n_j\), whose mean is \(\mu_j-n_j\le0\) by~\eqref{eq:policy_constraints}. If \(\mu_j=n_j\), the same random-walk maximum argument used in the single-type proof gives
\[
    \mathbb{E}[\bar W^B_{(j)}]
    \le \mathbb{E}[\widetilde W^B_{(j)}]
    =O(B^{1/2}).
\]
The same bound holds for every prefix \(b\le B\), so
\[
    \frac{1}{B}\sum_{b=0}^{B-1}\mathbb{E}[\bar W^b_{(j)}]=O(B^{1/2}).
\]
If \(\mu_j<n_j\), the dominating process has strictly negative drift. Since
\[
\left|\mathbb{E}[\bar X^b_{(j)}-n_j]\right|=n_j-\mu_j,
\]
by Kingman's bound \citep{kingman1962queues}:
\[
    \sup_{B\ge0}\mathbb{E}[\bar W^B_{(j)}]
    \le
    2n_j+
    \frac{\mathbb{E}[(\bar X^b_{(j)}-n_j)^2]}{2(n_j-\mu_j)}
    <\infty .
\]
Combining the zero-drift and negative-drift cases, under the nonpositive drift condition in~\eqref{eq:policy_constraints},
\[
    \sum_{j=1}^m \mathbb{E}[\bar W^B_{(j)}]=O(B^{1/2}),
    \qquad
    \frac{1}{B}\sum_{b=0}^{B-1}\sum_{j=1}^m \mathbb{E}[\bar W^b_{(j)}]=O(B^{1/2}).
\]
Under strict slack \(\mu_j<n_j\) for all \(j\), both quantities are \(O(1)\). If equality holds for at least one type and strict inequality holds for the others, the zero-drift types dominate the aggregate bounds, so the first part of the theorem applies. Types with zero arrival rates can be omitted from \([m]\).

It remains to translate these \(B\)-slot bounds into the theorem's finite-horizon metrics. The number of completed comparison slots in the physical horizon \([0,T]\) is exactly
\[
    B_\zeta(T)=\max\{b:u_b\le T\}
    =\left\lfloor\frac{\zeta T}{\Delta^\pi}\right\rfloor,
\]
and the residual time is less than one comparison slot. This is the number of deterministic comparison slots completed by \(\bar\pi\), not the number of actual event-driven WAIT batches. Since \(B_\zeta=\zeta T/\Delta^\pi+O(1)\), \(B_\zeta^{-1/2}=O((\zeta T)^{-1/2})\) and \(B_\zeta^{-1}=O((\zeta T)^{-1})\).

For throughput, terminal entry backlog controls the completion shortfall. More precisely, for the comparison policy,
\[
    \throughput^*-\mathbb{E}[\throughput^{(\zeta,\bar\pi)}]
    \le
    \frac{1}{\zeta T}\sum_{j=1}^m l_j^\prime\,
    \mathbb{E}\!\left[
    \bar W^{B_\zeta}_{(j)}
    +A_j(u_{B_\zeta},T]
    +(l_j^\prime+2)n_j
    \right].
\]
The second term accounts for final partial-slot arrivals, and the last term absorbs the bounded in-service pipeline, startup effects, and the one-slot comparison lag from Lemma~\ref{lem:wait_one_slot_dominance}. Since \(T-u_{B_\zeta}<\Delta^\pi/\zeta\), \(\mathbb{E}[A_j(u_{B_\zeta},T]]\le \lambda_j\Delta^\pi=O(1)\). Thus the terminal-backlog term gives an \(O((\zeta T)^{-1/2})\) gap under nonpositive drift and an \(O((\zeta T)^{-1})\) gap under strict slack.

For delay, terminal backlog is not the right object; the relevant quantity is the time average of the entry queues. Let \(\bar Q^b=\sum_{j=1}^m \bar W^b_{(j)}\). In physical time, each comparison slot has length \(\Delta^\pi/\zeta\). The theorem, however, reports service-normalized delay, multiplying elapsed physical time by \(\zeta\). Thus each comparison slot contributes a constant scaled time \(\Delta^\pi\). The entry-wait integral in service-normalized units is bounded by
\[
    \Delta^\pi\sum_{b=0}^{B_\zeta-1}\bar Q^b+O(B_\zeta),
\]
where the \(O(B_\zeta)\) term accounts for the one-slot convention and bounded resident-stage waiting. The expected number of completed prompts is \(\Theta(B_\zeta)\): cumulative cohort removals equal cumulative arrivals minus terminal backlog, and the missed removals are \(O(B_\zeta^{1/2})\) at the boundary and \(O(1)\) under strict slack. Therefore
\[
    \mathbb{E}[\Latency^{(\zeta,\bar\pi)}],
    \mathbb{E}[\TTFT^{(\zeta,\bar\pi)}]
    \le
    O\!\left(
    \frac{1}{B_\zeta}\sum_{b=0}^{B_\zeta-1}
    \mathbb{E}[\bar Q^b]
    \right)+O(1).
\]
The prefix bound gives
\[
    \mathbb{E}[\Latency^{(\zeta,\bar\pi)}],
    \mathbb{E}[\TTFT^{(\zeta,\bar\pi)}]
    =
    O(B_\zeta^{1/2})
    =
    O((\zeta T)^{1/2})
\]
under nonpositive drift. Under strict slack, the same expression is \(O(1)\). Lemma~\ref{lem:wait_one_slot_dominance} transfers these throughput, latency, and TTFT bounds from \(\bar\pi\) to the actual event-driven WAIT policy, completing the proof of Theorem~\ref{thm:wait_heavy_traffic}.
